% Plantilla para Trabajos Finales en la EPEF (UDC)
\documentclass{Plantilla/Formato_EPEF}

% Los paquetes están incluidos en la plantilla, no en este documento

%%%%%%%%%%%%%%%%%%%%%%%%%%%%%
% Datos relativos al alumno %
%%%%%%%%%%%%%%%%%%%%%%%%%%%%%
% Nombre del autor del trabajo
\newcommand{\Autor}{Nombre del autor del Trabajo}
% Sexo del autor (comentar el que no sea correcto)
\newcommand{\AutorSex}{Autor}
%\newcommand{\AutorSex}{Autora}
% Titulación para la que se presenta el Trabajo Final (comentar el que no sea correcto)
\newcommand{\Titulacion}{Elegir la titulación correspondiente}
% \newcommand{\Titulacion}{Grado Dual en Ingeniería Eléctrica}
% \newcommand{\Titulacion}{Grado en Ingeniería Electrónica Industrial y Automática}
% \newcommand{\Titulacion}{Grado Abierto en Ingeniería Industrial}
% \newcommand{\Titulacion}{Grado en Ingeniería en Tecnologías Industriales}
% \newcommand{\Titulacion}{Grado en Ingeniería Mecánica}
% \newcommand{\Titulacion}{Grado en Ingeniería Naval y Oceánica}
% \newcommand{\Titulacion}{M. U. en Eficiencia Energética y Sostenibilidad}
% \newcommand{\Titulacion}{M. U. en Informática Industrial y Robótica}
% \newcommand{\Titulacion}{M. U. en Diseño, Desarrollo y Comercialización de Videojuegos}
% \newcommand{\Titulacion}{M. U. en Ingeniería Industrial}
% \newcommand{\Titulacion}{M. U. en Ingeniería Naval y Oceánica}
% \newcommand{\Titulacion}{M. U. en Materiales Complejos: Análisis Térmico y Reología}
% \newcommand{\Titulacion}{M. U. en Erasmus Mundus en Sostenibilidad e Industria 4.0 aplicada al Sector Marítimo}
% Nombre de la escuela
\newcommand{\Escuela}{Escuela Politécnica de Ingeniería de Ferrol}
\newcommand{\EscuelaAbrevidado}{EPEF}
% Dirección de la escuela
\newcommand{\Direccion}{Rúa Mendizábal, s/n, Campus de Esteiro, 15403, Ferrol}

%%%%%%%%%%%%%%%%%%%%%%%%%%%%%%%%%%%%
% Datos relativos al Trabajo Final %
%%%%%%%%%%%%%%%%%%%%%%%%%%%%%%%%%%%%
% Título del Trabajo Final
\newcommand{\TituloUno}{Título del Trabajo Final}
\newcommand{\TituloDos}{} % Usar sólo si el título es demasiado largo
\newcommand{\TituloTres}{} % Usar sólo si el título es demasiado largo
% Número asociado al Trabajo Final
\newcommand{\Numero}{xxx}
% Tipo de Trabajo Final (comentar el que no sea correcto)
%\newcommand{\TipoTrabajo}{Trabajo Fin de Grado}
%\newcommand{\TipoAbreviado}{TFG}
\newcommand{\TipoTrabajo}{Trabajo Fin de Máster}
\newcommand{\TipoAbreviado}{TFM}
% Tutor/es (en caso de no haber co-tutor, se deberá de comentar la línea correspondiente)
\newcommand{\TutorUno}{Nombre del tutor}
\newcommand{\TutorDos}{} % Usar sólo en caso de haber un co-tutor
% Sexo y número del tutor (comentar el que no sea correcto)
\newcommand{\TutorSex}{Tutor}
%\newcommand{\TutorSex}{Tutora}
%\newcommand{\TutorSex}{Tutores}
%\newcommand{\TutorSex}{Tutoras}
% Fecha de presentación del Trabajo Final
\newcommand{\Fecha}{Noviembre de 2022}
\newcommand{\Curso}{2022/2023}

\begin{document}
\subfile{Plantilla/Portada_principal}

\tableofcontents
\newpage
\listoffigures
\newpage
\listoftables
\cleardoublepage

%%%%%%%%%%%%%%%%%%%%%%%%%%%%%%%%%%
% Resúmenes en todos los idiomas %
%%%%%%%%%%%%%%%%%%%%%%%%%%%%%%%%%%
\addcontentsline{toc}{section}{Resumen}
\section*{Resumen}
Resumen del trabajo desarrollado.

\section*{Resumo}
Resumo do traballo en galego.

\section*{Abstract}
Abstract in English.

%%%%%%%%%%%%%%%%%%%%%%%%%%%%%%%%%%%
% Contenido principal del trabajo %
%%%%%%%%%%%%%%%%%%%%%%%%%%%%%%%%%%%
\newpage
\section{Introducción}
Contextualización del TFM desarrollado, describir los antecedentes y estado del arte que apliquen al mismo.

\newpage
\section{Objetivos}
Describir los objetivos del trabajo, así como el grado de consecución de los mismos.

\newpage
\section{Materiales y métodos}
Descripción de las tecnologías, herramientas, técnicas y conjuntos de datos.

\newpage
\section{Experimentos / Desarrollo}
Descripción del trabajo, así como los experimentos realizados.

\newpage
\section{Resultados}
Presentación de los resultados obtenidos en forma de tablas, gráficos, figuras, etc. y análisis de los mismos en base a datos objetivos presentados.

\newpage
\section{Discusión}
(Sólo si es de aplicación en  base a la temática del trabajo) Comparación de los resultados obtenidos con investigación previas, mostrando los avances científicos obtenidos en relación al estado actual de las técnicas. Ejemplo de cita bibliográfica \cite{ref:Hydrogen}.

\newpage
\section{Conclusiones y trabajos futuros}
Presentación de las conclusiones y exposición de los trabajos futuros a los que puede dar lugar el trabajo.

\newpage
\bibliographystyle{plain}
\bibliography{Bibliography}

\newpage
\section{Documentación de partida}
Documento de asignación del TFM donde se refleja la descripción del trabajo a realizar y los objetivos planteados.

\end{document}
